% Use the existing owner-derived source mechanism; do not edit frozen originals.
\OLSADeclareDocumentTextCorrection{OLP-0031}{FA0031-DIRECT-INVERSE-EXPOSITION}
{خوب بود اگر می‌توانستیم
این مقدار را مستقیم‌تر محاسبه کنیم. یعنی خوب بود \emph{وارونِ} برشماری
زیگزاگ، یعنی $g\colon \Nat^2 \to \Nat$، را در دسترس داشتیم، به‌گونه‌ای که}
{اگر بتوانیم
این مقدار را مستقیماً محاسبه کنیم، کار ساده‌تر می‌شود. به بیان دیگر، می‌خواهیم
\emph{معکوسِ} برشماری زیگزاگ، یعنی $g\colon \Nat^2 \to \Nat$، را به‌صورت صریح
داشته باشیم؛ به‌گونه‌ای که}
\OLSADeclareDocumentTextCorrection{OLP-0031}{FA0031-PRESENT-FUNCTION-PROPERTY}
{این تابع به ما امکان می‌داد}
{این تابع به ما امکان می‌دهد}
\OLSADeclareDocumentTextCorrection{OLP-0031}{FA0031-POSITIVE-COLUMN}
{حاوی مقدار~$v$ باشد، آنگاه سطر}
{حاوی مقدار~$v$ باشد و شمارهٔ ستون مثبت باشد، آنگاه سطر}
\OLSADeclareDocumentTextCorrection{OLP-0031}{FA0031-TRIANGULAR-SUM}
{اعداد طبیعیِ $< k$}
{اعداد طبیعیِ $\leq k$}
\OLSADeclareDocumentTextCorrection{OLP-0031}{FA0031-SOURCE-BOUNDARY-DISCLOSURE}
{محاسبه کرد. با کنار هم گذاشتن این}
{محاسبه کرد. \emph{یادداشت ویراستاری:} شرط مثبت‌بودن شمارهٔ ستون در مشاهدهٔ نخست،
و احتساب خودِ عدد در مجموع اعداد مثلثی، دو اصلاح در توضیح متن مبدأ هستند.
با کنار هم گذاشتن این}
\OLSADeclareDocumentTextCorrection{OLP-0031}{FA0031-INVERSE-TERM}
{این تابع $g$ \emph{وارونِ}}
{این تابع $g$ \emph{معکوسِ}}
\OLSADeclareDocumentTextCorrection{OLP-0031}{FA0031-DECODING-DOMAIN}
{با استفاده از وارون تابع جفت‌سازی می‌توانیم
عدد را \emph{کدگشایی کنیم}؛ یعنی دریابیم کدام زوج را بازنمایی می‌کند.}
{با استفاده از معکوس تابع جفت‌سازی، که بر برد آن تعریف شده است، می‌توانیم
\emph{کدگشایی کنیم}؛ یعنی زوجی را که یک کدِ معتبر بازنمایی می‌کند به دست آوریم.
\emph{توضیح ویراستاری:} یک‌به‌یک‌بودن به‌تنهایی تضمین نمی‌کند که هر عددی کدِ یک زوج باشد.}
\OLSADeclareDocumentTextCorrection{OLP-0031}{FA0031-TRUTH-FUNCTION-SCOPE}
{از درس مقدماتی منطق خود به یاد آورید که هر جدول ارزشِ ممکن یک تابع ارزش
را بیان می‌کند. به بیان دیگر، توابع ارزش همهٔ توابع از
$\Bin^k \to \Bin$ برای یک~$k$ هستند. ثابت کنید مجموعهٔ همهٔ توابع ارزش
شمارا است.}
{از درس مقدماتی منطق خود به یاد آورید که هر جدول ارزشِ ممکن یک تابعِ ارزشِ صدق
را تعیین می‌کند. به بیان دیگر، این‌ها همهٔ توابعی به‌صورت
$\Bin^k \to \Bin$ برای یک~$k$ هستند. ثابت کنید مجموعهٔ همهٔ توابعِ ارزشِ صدق
شماراست.}
\OLSADeclareDocumentTextCorrection{OLP-0031}{FA0031-COFINITE-DISCLOSURE}
{نشان دهید $I$ !!{enumerable} است.}
{نشان دهید $I$ !!{enumerable} است.
\emph{یادداشت ویراستاری:} عبارت ناقصِ متن مبدأ دربارهٔ مجموعهٔ متناهی، مطابق
تعریف فرمولیِ بلافاصله پس از آن، «زیرمجموعهٔ متناهیِ اعداد طبیعی» خوانده شده است.}
\OLSADeclareDocumentTextCorrection{OLP-0031}{FA0031-COUNTABLE-UNION-CHOICE}
{[توجه: این مسئله دشوار است!]}
{[توجه: این مسئله دشوار است!]
\emph{یادداشت دربارهٔ فرض‌ها:} در برهان معمولِ این گزاره، برای انتخاب هم‌زمانِ
یک برشماری برای هر مجموعهٔ ناتهیِ خانواده از اصل انتخاب استفاده می‌شود.
اگر این برشماری‌ها از ابتدا جزو داده‌های مسئله باشند، می‌توان همان داده‌ها را
در ساخت قطری به کار برد؛ در آن صورت این مرحلهٔ انتخاب لازم نیست.}
\OLSADeclareDocumentTextCorrection{OLP-0031}{FA0031-INVERSE-EXERCISE}
{نشان دهید وارونِ $f$ یک برشماری از $A \times B$ است.}
{\emph{صورت اصلاح‌شدهٔ تمرین:} در حالت ناتهی، با استفاده از معکوسِ $f$
که بر برد آن تعریف شده است، برشماری‌ای برای $A \times B$ بسازید.
حالت تهی را نیز جداگانه بررسی کنید.
\emph{یادداشت ویراستاری:} صورت اولیهٔ تمرین، معکوس را بی‌قید برشماری می‌نامید؛
اما معکوسِ یک تابع یک‌به‌یک لزوماً بر همهٔ اعداد طبیعی تعریف نشده است.}
\OLSADeclareDocumentTextCorrection{OLP-0031}{FA0031-CELL-NOT-ROW}
{سطر $n$اُم و ستون $m$اُم حاوی مقدار~$v$ باشد و شمارهٔ ستون مثبت باشد، آنگاه سطر
$(n+1)$اُم و ستون $(m-1)$اُم حاوی}
{خانهٔ واقع در سطر $n$اُم و ستون $m$اُم حاوی مقدار~$v$ باشد و شمارهٔ ستون مثبت باشد،
آنگاه خانهٔ واقع در سطر $(n+1)$اُم و ستون $(m-1)$اُم حاوی}
\OLSADeclareDocumentTextCorrection{OLP-0032}{FA0032-INVERSE-TERM}
{یافتن وارونشان}
{یافتن معکوسشان}
\OLSADeclareDocumentTextCorrection{OLP-0032}{FA0032-SECOND-PAIR}
{دومین زوج (یعنی $\tuple{0,2}$)}
{دومین زوج (یعنی $\tuple{0,1}$)}
\OLSADeclareDocumentTextCorrection{OLP-0032}{FA0032-SUCCESSIVE-ROWS}
{زوج‌های $\tuple{2,m}$، $\tuple{2,m}$ و جز آن}
{زوج‌های $\tuple{2,m}$، $\tuple{3,m}$ و جز آن}
\OLSADeclareDocumentTextCorrection{OLP-0032}{FA0032-TYPO-DISCLOSURE}
{کنید. بنابراین برشماری کامل‌شدهٔ ما چنین آغاز می‌شود:}
{کنید. \emph{یادداشت ویراستاری:} دو لغزش در مؤلفه‌های زوج‌ها در توضیح متن مبدأ
اصلاح شده‌اند تا توضیح با جدول‌ها سازگار باشد.
بنابراین برشماری کامل‌شدهٔ ما چنین آغاز می‌شود:}
\OLSADeclareDocumentTextCorrection{OLP-0032}{FA0032-ODD-POSITION-IDIOM}
{در جایگاه‌های فردشمارهٔ برشماری ما
قرار دارند}
{در جایگاه‌هایی از برشماری ما قرار دارند
که شمارهٔ آن‌ها فرد است}
\OLSADeclareDocumentTextCorrection{OLP-0032}{FA0032-MULTIPLIER-IDIOM}
{عامل‌های
$(2m+1)$ برای هر سطر}
{ضریب‌های
$(2m+1)$ در سطرهای مختلف}
\OLSADeclareDocumentTextCorrection{OLP-0032}{FA0032-EXPONENT-IDIOM}
{نمای مربوط همواره مؤلفهٔ نخست زوج مورد بحث است.}
{این نما همواره مؤلفهٔ نخستِ زوج مورد نظر است.}
\OLSADeclareDocumentTextCorrection{OLP-0032}{FA0032-MULTIPLIER-SINGLE}
{عامل برابر $2^n$}
{ضریب برابر $2^n$}
\OLSADeclareDocumentTextCorrection{OLP-0032}{FA0032-PARTIAL-INVERSE}
{لازم باشد کدگذاری پوشا باشد. چنین کدگذاری‌هایی وارون‌هایی دارند که فقط
توابع جزئی‌اند.}
{لازم باشد کدگذاری پوشا باشد. اگر کدگذاری پوشا نباشد، معکوس آن فقط در بردش
تعریف می‌شود؛ بنابراین با در نظر گرفتن همهٔ اعدادِ مجموعهٔ مقصد به‌عنوان
ورودی‌های ممکن، معکوس یک تابع جزئی است.}
\OLSADeclareSetup{OLP-0031}{\olsa_source_correction_install_document_text:}
\OLSADeclareSetup{OLP-0032}{\olsa_source_correction_install_document_text:}
